% ==============================================================================
% CHƯƠNG 7: TRA CỨU THUỘC TÍNH FORM & CÁC LOẠI INPUT TYPES
% ==============================================================================
\chapter{Tra Cứu Thuộc Tính Form \& Các Loại Input Types}
\label{chap:forms_masterclass}

Chương này đi sâu phân tích toàn bộ các thẻ biểu mẫu HTML Forms, bao gồm 22+ loại Input Types, các thuộc tính điều khiển và cơ chế kiểm tra hợp lệ Constraint Validation API.

\section{Thẻ Khung Biểu Mẫu <form>}
\subsection{Lý Thuyết \& Bảng Toàn Bộ Thuộc Tính}
Thẻ \codeinline{<form>} đại diện cho một phần tử chứa các trường nhập liệu tương tác để thu thập và gửi dữ liệu tới máy chủ.

\begin{table}[H]
\centering
\begin{prettyTableBox}
\small
\rowcolors{2}{white}{primaryBlue!4}
\begin{tabularx}{\linewidth}{lp{3.5cm}X}
\rowcolor{primaryBlue}
\color{white}\textbf{Thuộc Tính} & \color{white}\textbf{Giá Trị Hợp Lệ} & \color{white}\textbf{Chức Năng \& Hành Vi} \\
\hline
\codeinline{action} & URL máy chủ & Địa chỉ API nhận dữ liệu form khi submit. \\
\codeinline{method} & \codeinline{GET}, \codeinline{POST}, \codeinline{dialog} & Phương thức truyền dữ liệu HTTP Request. \\
\codeinline{enctype} & \codeinline{application/x-www-form-urlencoded}, \codeinline{multipart/form-data} & Mã hóa dữ liệu (Dùng \codeinline{multipart} khi tải file). \\
\codeinline{autocomplete} & \codeinline{on}, \codeinline{off} & Bật/Tắt tự động gợi ý điền thông tin của trình duyệt. \\
\codeinline{novalidate} & Boolean & Tắt cơ chế kiểm tra hợp lệ dữ liệu mặc định của HTML5. \\
\codeinline{target} & \codeinline{\_self}, \codeinline{\_blank} & Nơi nhận kết quả trả về từ server. \\
\end{tabularx}
\end{prettyTableBox}
\vspace{-2pt}
\caption{Toàn bộ các thuộc tính của thẻ <form>}
\end{table}

\subsection{Ví Dụ Mã Nguồn Thực Tế}
\begin{lstlisting}[language=HTML5, caption={Ví dụ biểu mẫu đăng ký hồ sơ hỗ trợ tải file}]
<form action="/api/v1/register" method="POST"
      enctype="multipart/form-data" autocomplete="on">
  <!-- (*@Các trường nhập liệu nằm ở đây@*) -->
  <button type="submit">(*@Đăng Ký Hồ Sơ@*)</button>
</form>
\end{lstlisting}

% ------------------------------------------------------------------------------
\section{Thẻ Nhãn Dán <label>}
\subsection{Lý Thuyết \& Accessibility (a11y)}
Thẻ \codeinline{<label>} tạo nhãn mô tả cho ô nhập liệu. Khi người dùng nhấp chuột vào thẻ \codeinline{<label>}, trình duyệt sẽ tự động nhảy con trỏ chuột vào ô input tương ứng.

\begin{itemize}
    \item Thuộc tính \codeinline{for="id\_input"}: Liên kết nhãn với \codeinline{id} của phần tử input.
\end{itemize}

\subsection{Ví Dụ Mã Nguồn Thực Tế}
\begin{lstlisting}[language=HTML5, caption={Ví dụ liên kết thẻ label với input}]
<label for="username">(*@Tên Đăng Nhập@*)<span style="color: red;">*</span></label>
<input type="text" id="username" name="username" required>
\end{lstlisting}

% ------------------------------------------------------------------------------
\section{Thẻ Khung Nhóm <fieldset> \& Tiêu Đề Khung <legend>}
\subsection{Lý Thuyết \& Bảng Thuộc Tính}
- \codeinline{<fieldset>}: Gom nhóm các ô nhập liệu có liên quan lại trong một khung viền.
- \codeinline{<legend>}: Đặt tiêu đề hiển thị nằm trên viền của khối \codeinline{<fieldset>}.
- Thuộc tính \codeinline{disabled} trên \codeinline{<fieldset>}: Vô hiệu hóa \textbf{tất cả các input con} bên trong cùng lúc!

\subsection{Ví Dụ Mã Nguồn Thực Tế}
\begin{lstlisting}[language=HTML5, caption={Ví dụ gom nhóm trường thông tin bằng fieldset và legend}]
<fieldset>
  <legend>(*@Thông Tin Giao Hàng@*)</legend>
  <label for="address">(*@Địa Chỉ:@*)</label>
  <input type="text" id="address" name="address">
</fieldset>
\end{lstlisting}

% ------------------------------------------------------------------------------
\section{Thẻ Đầu Vào <input> \& Toàn Bộ 22+ Input Types}
\subsection{Lý Thuyết \& Bảng Tra Cứu All Input Types}
Thẻ \codeinline{<input>} là thẻ nhập liệu quan trọng nhất trong HTML5, linh hoạt thay đổi giao diện dựa vào thuộc tính \codeinline{type="..."}.

\begin{table}[H]
\centering
\begin{prettyTableBox}
\small
\rowcolors{2}{white}{primaryBlue!4}
\begin{tabularx}{\linewidth}{lX}
\rowcolor{primaryBlue}
\color{white}\textbf{Input Type} & \color{white}\textbf{Giao Diện \& Hành Vi Trình Duyệt} \\
\hline
\codeinline{text}, \codeinline{password} & Nhập văn bản 1 dòng và ẩn mật khẩu dạng dấu chấm. \\
\codeinline{email}, \codeinline{url}, \codeinline{tel} & Tự động bật bàn phím di động tương ứng (@, .com, phím số). \\
\codeinline{search} & Ô tìm kiếm có nút dấu "X" xóa nhanh văn bản. \\
\codeinline{number}, \codeinline{range} & Nhập số có nút tăng giảm và thanh trượt khoảng. \\
\codeinline{color} & Bật bảng chọn màu sắc (Color Picker). \\
\codeinline{date}, \codeinline{time}, \codeinline{datetime-local} & Bật giao diện chọn Lịch (Calendar) và Đồng hồ gốc. \\
\codeinline{month}, \codeinline{week} & Chọn tháng/năm hoặc tuần cụ thể. \\
\codeinline{checkbox}, \codeinline{radio} & Chọn nhiều ô vuông hoặc chọn duy nhất 1 ô tròn trong nhóm \codeinline{name}. \\
\codeinline{file} & Tải tệp lên máy chủ (Hỗ trợ thuộc tính \codeinline{multiple}, \codeinline{accept}, \codeinline{capture}). \\
\codeinline{hidden} & Trữ dữ liệu ngầm gửi lên server mà không hiển thị ra màn hình. \\
\codeinline{submit}, \codeinline{reset}, \codeinline{button} & Các dạng nút thực thi form cổ điển. \\
\end{tabularx}
\end{prettyTableBox}
\vspace{-2pt}
\caption{Bảng 22+ loại kiểu đầu vào của thẻ <input>}
\end{table}

\subsection{Bảng Thuộc Tính Nhập Liệu \& Mobile UX Attributes}
\begin{itemize}
    \item \codeinline{inputmode="numeric|decimal|tel|search|email"}: Ép bàn phím ảo hiển thị đúng dạng.
    \item \codeinline{enterkeyhint="search|send|done|next"}: Đổi nhãn nút Enter trên bàn phím di động.
    \item \codeinline{autocomplete="one-time-code"}: Tự động đọc mã OTP SMS gửi về di động.
    \item \codeinline{pattern="..."}: Ràng buộc dữ liệu bằng Biểu thức chính quy (Regex).
\end{itemize}

\subsection{Ví Dụ Mã Nguồn Thực Tế}
\begin{lstlisting}[language=HTML5, caption={Ví dụ tổng hợp các ô input di động nâng cao}]
<!-- (*@Ô nhập OTP tự động đọc SMS trên iOS/Android@*) -->
<label for="otp">(*@Nhập Mã OTP:@*)</label>
<input type="text" id="otp" name="otp" inputmode="numeric" 
       autocomplete="one-time-code" required placeholder="(*@6 chữ số@*)">

<!-- (*@Chọn tệp ảnh PNG/JPEG nhiều file@*) -->
<label for="photos">(*@Tải Ảnh Hóa Đơn:@*)</label>
<input type="file" id="photos" name="photos" 
       accept="image/png, image/jpeg" multiple>

<!-- (*@Thanh chọn khoảng giá@*) -->
<label for="price">(*@Khoảng Giá:@*)</label>
<input type="range" id="price" name="price" 
       min="100" max="1000" step="50" value="500">
\end{lstlisting}

% ------------------------------------------------------------------------------
\section{Thẻ Nút Bấm <button>}
\subsection{Lý Thuyết \& Các Thuộc Tính Ghi Đè Form}
Thẻ \codeinline{<button>} tạo nút bấm tương tác linh hoạt hơn thẻ \codeinline{<input type="button">} vì cho phép nhúng icon HTML bên trong.

\begin{itemize}
    \item \codeinline{type="submit|reset|button"}: Loại hành động của nút (Mặc định là \codeinline{submit}).
    \item Thuộc tính ghi đè Form: \codeinline{formaction}, \codeinline{formenctype}, \codeinline{formmethod}, \codeinline{formnovalidate}, \codeinline{formtarget}.
\end{itemize}

\subsection{Ví Dụ Mã Nguồn Thực Tế}
\begin{lstlisting}[language=HTML5, caption={Ví dụ nút bấm button chứa icon SVG}]
<button type="submit" class="btn-primary">
  <svg width="16" height="16"><circle cx="8" cy="8" r="7"/></svg>
  <span>(*@Gửi Đơn Hàng@*)</span>
</button>
\end{lstlisting}

% ------------------------------------------------------------------------------
\section{Thẻ Danh Sách Chọn <select>, <optgroup> \& <option>}
\subsection{Lý Thuyết \& Bảng Thuộc Tính}
- \codeinline{<select>}: Thẻ tạo danh sách chọn xổ xuống (Dropdown). Thuộc tính \codeinline{multiple} (cho chọn nhiều), \codeinline{size} (số dòng hiển thị).
- \codeinline{<optgroup>}: Gom nhóm các tùy chọn có cùng chủ đề. Thuộc tính \codeinline{label}.
- \codeinline{<option>}: Một phần tử chọn cụ thể. Thuộc tính \codeinline{value}, \codeinline{selected}, \codeinline{disabled}.

\subsection{Ví Dụ Mã Nguồn Thực Tế}
\begin{codeWindow}[Mã Nguồn: Menu Dropdown Gom Nhóm Optgroup]
\begin{lstlisting}[language=HTML5]
<label for="city">(*@Chọn Thành Phố:@*)</label>
<select id="city" name="city">
  <optgroup label="(*@Miền Bắc@*)">
    <option value="hn" selected>(*@Hà Nội@*)</option>
    <option value="hp">(*@Hải Phòng@*)</option>
  </optgroup>
  <optgroup label="(*@Miền Nam@*)">
    <option value="hcm">(*@Hồ Chí Minh@*)</option>
    <option value="ct">(*@Cần Thơ@*)</option>
  </optgroup>
</select>
\end{lstlisting}
\end{codeWindow}

\begin{browserWindow}[Kết Quả Hiển Thị Trình Duyệt: Select Gom Nhóm Thành Phố Optgroup]
\sffamily\small
\textbf{Chọn Thành Phố:} \par\vspace{4pt}
\begin{tcolorbox}[arc=2pt, colback=white, colframe=primaryBlue!50, boxsep=4pt, width=0.55\linewidth]
  \small
  \textbf{\textcolor{gray!80}{MIỀN BẮC}} \par\vspace{2pt}
  \quad \textcolor{primaryBlue}{\textbf{Hà Nội (Được chọn sẵn)}} \par\vspace{2pt}
  \quad Hải Phòng \par\vspace{3pt}
  \textbf{\textcolor{gray!80}{MIỀN NAM}} \par\vspace{2pt}
  \quad Hồ Chí Minh \par\vspace{2pt}
  \quad Cần Thơ
\end{tcolorbox}
\par\vspace{4pt}
\textcolor{codeComment}{\small\textit{$\rightarrow$ Thẻ optgroup phân chia Miền Bắc / Miền Nam in đậm, mục "Hà Nội" được chọn mặc định nhờ thuộc tính selected.}}
\end{browserWindow}

% ------------------------------------------------------------------------------
\section{Thẻ Ô Nhập Nhiều Dòng <textarea>}
\subsection{Lý Thuyết \& Thuộc Tính Khoảng Kích Thước}
Thẻ \codeinline{<textarea>} cho phép người dùng nhập các đoạn văn bản dài nhiều dòng (như bình luận, ghi chú đơn hàng).

\begin{itemize}
    \item \codeinline{rows="..."}: Số dòng hiển thị mặc định.
    \item \codeinline{cols="..."}: Số cột chiều rộng mặc định.
    \item \codeinline{maxlength="..."}: Giới hạn số ký tự tối đa.
\end{itemize}

\subsection{Ví Dụ Mã Nguồn Thực Tế}
\begin{lstlisting}[language=HTML5, caption={Ví dụ khung textarea nhập bình luận}]
<label for="comment">(*@Bình Luận Bài Viết:@*)</label>
<textarea id="comment" name="comment" rows="5" cols="50" 
          maxlength="500" placeholder="(*@Viết phản hồi của bạn...@*)"></textarea>
\end{lstlisting}

% ------------------------------------------------------------------------------
\section{Thẻ Gợi Ý Tự Động <datalist>}
\subsection{Lý Thuyết \& Cơ Chế Tự Động Điền}
Thẻ \codeinline{<datalist>} chứa danh sách các tùy chọn gợi ý cho thẻ \codeinline{<input>}. Người dùng vừa có thể gõ văn bản tùy ý, vừa có thể chọn từ danh sách gợi ý.

\subsection{Ví Dụ Mã Nguồn Thực Tế}
\begin{lstlisting}[language=HTML5, caption={Ví dụ liên kết ô input với thẻ datalist}]
<label for="browser">(*@Chọn Trình Duyệt Thường Dùng:@*)</label>
<input type="text" id="browser" name="browser" list="browser-list">

<datalist id="browser-list">
  <option value="Google Chrome">
  <option value="Mozilla Firefox">
  <option value="Apple Safari">
  <option value="Microsoft Edge">
</datalist>
\end{lstlisting}

% ------------------------------------------------------------------------------
\section{Thẻ Hiển Thị Kết Quả <output>}
\subsection{Lý Thuyết \& Thuộc Tính for}
Thẻ \codeinline{<output>} đại diện cho kết quả của một phép tính toán hoặc hành động người dùng thực hiện trong form.

\begin{itemize}
    \item Thuộc tính \codeinline{for="id1 id2"}: Khai báo danh sách ID các ô input tham gia vào phép tính.
\end{itemize}

\subsection{Ví Dụ Mã Nguồn Thực Tế}
\begin{lstlisting}[language=HTML5, caption={Ví dụ tự động tính toán tổng số tiền bằng thẻ output và JS}]
<form oninput="result.value = parseInt(a.value) * parseInt(b.value)">
  <input type="number" id="a" value="10"> x
  <input type="number" id="b" value="5"> =
  (*@<output name="result" for="a b">50</output> VNĐ@*)
</form>
\end{lstlisting}

\section{Tóm Tắt Chương 7}
Chương 7 đã hệ thống trọn bộ các mục cho từng thẻ Biểu mẫu HTML Forms, 22+ loại Input Types, thuộc tính bàn phím di động và thẻ tính toán kết quả \codeinline{<output>}. Chương tiếp theo sẽ đi sâu vào nhóm Thẻ Đa Phương Tiện và Đồ Họa Graphics.
